\chapter{Distribution des données de marché et droits d'accès}\label{ch:pl:market-data-distribution-and-entitlements}

L'auditeur d'une bourse demande à une société la liste de toutes les personnes et de tous les programmes qui ont utilisé les données de profondeur de la bourse au cours des trois dernières années.
La société peut citer quarante écrans~: elle en paie quarante. Elle ne peut pas citer, dans un premier temps, les onze applications qui lisent le même flux depuis la
plateforme de distribution de la société --- stratégies, un service de risque, un enregistreur de ticks, un système de surveillance, un tableau de bord de recherche --- et chacune d'elles représente
un usage facturé par la bourse. La société de ce chapitre est un modèle, mais les documents selon lesquels elle est auditée sont réels~: la grille tarifaire
d'une bourse, ses politiques de données et son contrat d'abonnement. Le chapitre construit la plateforme qui distribue les données de marché au sein d'une société, et la partie de celle-ci
qui fait le plus souvent défaut~: le registre de qui a utilisé quoi, dans quel but, compté de la façon dont la plateforme le compte.

\section{Distribuer les données au sein de la société}

Un flux direct (livre~1, chapitre~9) arrive une seule fois aux feed handlers de la société. De là, il doit atteindre des dizaines de consommateurs aux besoins très différents
besoins~: une stratégie qui veut chaque message en quelques microsecondes, un service de risque qui veut les cours courants, l'écran d'un trader qui ne peut afficher
plus de quelques mises à jour par seconde, un traitement de recherche qui enregistre tout. Relier chacun d'eux au flux multiplierait
la charge du réseau, les licences et les points de défaillance~; la société interpose une seule plateforme entre le flux et les consommateurs.

\begin{definition}[Plateforme de distribution des données de marché]\label{def:pl:market-data-distribution-and-entitlements:platform}
Une \emph{plateforme de distribution des données de marché}\index{plateforme de distribution des données de marché} reçoit une fois les données de marché normalisées et les diffuse
au sein de la société aux abonnés par thème (un instrument, une plateforme, un type de donnée), avec une politique de livraison propre à chaque abonné --- chaque message dans
l'ordre, ou en conflation --- et avec un contrôle des \omterm{def:pl:market-data-distribution-and-entitlements:entitlement}{droits d'accès} à chaque abonnement.
\end{definition}

\begin{omfigure}
\begin{tikzpicture}[font=\small,node distance=5mm,
  box/.style={draw=omDef,fill=omDef!8,rounded corners=2pt,align=center,minimum height=9mm,text width=22mm},
  chk/.style={draw=omThm,fill=omThm!8,rounded corners=2pt,align=center,minimum height=9mm,text width=22mm},
  arr/.style={-{Stealth[length=2mm]},thick}]
\node[box] (f) {feed handler\\(normalisé)};
\node[box,right=8mm of f,minimum height=34mm,text width=18mm] (b) {bus~:\\thèmes,\\une file par\\abonné};
\node[box,right=12mm of b,yshift=19.5mm,text width=36mm,minimum height=7mm] (s1) {stratégie, chaque message};
\node[box,right=12mm of b,yshift=6.5mm,text width=36mm,minimum height=7mm] (s2) {service de risque};
\node[box,right=12mm of b,yshift=-6.5mm,text width=36mm,minimum height=7mm] (s3) {tableau de bord de recherche};
\node[box,right=12mm of b,yshift=-19.5mm,text width=36mm,minimum height=7mm] (s4) {écran du trader (en conflation)};
\node[chk,below=12mm of f] (e) {système de\\droits d'accès};
\node[chk,right=4mm of e,xshift=22mm,yshift=-12mm] (u) {journal d'utilisation};
\node[chk,right=6mm of u] (r) {rapport d'utilisation\\et audit};
\draw[arr] (f) -- (b);
\foreach \x in {s1,s2,s3,s4} \draw[arr] (b.east) -- (\x.west);
\draw[arr,omThm] (e.east) -- node[above,font=\scriptsize,sloped]{contrôle} (b.south west);
\draw[arr,omThm] (b.south) |- (u.west);
\draw[arr,omThm] (u) -- (r);
\end{tikzpicture}

\omcaption{La plateforme de distribution du chapitre~: un flux normalisé, un bus avec une file par abonné, un contrôle des \omterm{def:pl:market-data-distribution-and-entitlements:entitlement}{droits d'accès} à l'abonnement
et à la livraison, et chaque accès inscrit dans un journal d'utilisation, à partir duquel sont
produits le \omterm{def:pl:market-data-distribution-and-entitlements:report}{rapport d'utilisation} mensuel et le \omterm{def:pl:post-trade-systems:recon}{rapprochement} de l'audit.}\label{fig:pl:market-data-distribution-and-entitlements:platform}
\end{omfigure}

Le bus du chapitre (\texttt{firm.entitle}) est simulé en temps virtuel. Dix minutes de mises à jour de profondeur pour vingt instruments --- 240\,267
messages, issus du générateur de tape du livre~7, le flux d'ordres que rejoue le simulateur du livre~10 --- sont destinés à quatre abonnés~: une stratégie qui prend
20~microsecondes par message, un service de risque qui prend 2~millisecondes, un tableau de bord de recherche qui prend 10, et l'écran d'un trader qui suit
cinq instruments à 50~millisecondes par rafraîchissement. Chaque abonné a sa propre file, de sorte qu'un abonné lent ne puisse pas retarder les autres~: la
conception publication--abonnement du chapitre~\ref{ch:pl:messaging-and-event-driven-architecture}, et le problème du consommateur lent
du livre~13, chapitre~12, transposé un niveau plus haut.

\omcode{code/firm/entitle/firm_entitle.py}{98}{111}{Les deux politiques de livraison~: une file de chaque mise à jour, dans l'ordre, ou la conflation, qui ne conserve que
la dernière mise à jour par thème et sert le thème qui attend depuis le plus longtemps.}\label{lst:pl:market-data-distribution-and-entitlements:queue}

\begin{table}[htbp]
\centering\small
\begin{tabular}{llrrrr}
\toprule
politique & abonné & livrés & file maximale & âge moyen & âge maximal\\
\midrule
chaque message & stratégie & 240\,267 & 1\,199 & 0.1~ms & 24~ms\\
 & service de risque & 240\,267 & 1\,203 & 58~ms & 2.4~s\\
 & tableau de bord de recherche & 240\,267 & 180\,267 & 891~s & 1\,803~s\\
conflation & stratégie & 238\,870 & 20 & 0.0~ms & 0.4~ms\\
 & service de risque & 221\,364 & 20 & 4.0~ms & 32~ms\\
 & tableau de bord de recherche & 60\,012 & 20 & 67~ms & 134~ms\\
\bottomrule
\end{tabular}
\caption{Dix minutes de mises à jour de profondeur pour vingt instruments livrés à trois abonnés, chaque message étant mis en file ou en conflation par
instrument. L'âge d'un message est le temps écoulé entre sa publication et la fin de son traitement. L'écran, en conflation dans les deux exécutions, voit un
âge moyen de 125~ms. La plus longue file de la stratégie est l'instantané d'ouverture, qui arrive en un seul instant.}
\label{tab:pl:market-data-distribution-and-entitlements:policies}
\end{table}

Le tableau de bord de recherche illustre la leçon du \cref{tab:pl:market-data-distribution-and-entitlements:policies}. Il peut traiter 100 messages par
seconde, et le flux en apporte environ 400~: avec chaque message mis en file, sa file grossit tout au long de la séance jusqu'à 180\,267 messages, et il affiche des cours
vieux de quinze minutes en moyenne, et il est encore en train de traiter la séance une demi-heure après sa fin
(\cref{fig:pl:market-data-distribution-and-entitlements:age}). En conflation, il traite un quart des messages et n'affiche jamais un cours
plus ancien que 134~ms. Le service de risque suit en moyenne mais pas lors des rafales~; la conflation ramène son âge maximal de 2.4~secondes à 32~ms. Un
écran ou un service de risque veut l'état courant, et non l'historique, et doit être en conflation~; une stratégie ou un enregistreur a besoin de chaque message
et doit être assez rapide pour s'en passer.

\begin{omfigure}
\begin{tikzpicture}
\begin{semilogyaxis}[width=11.5cm,height=5.6cm,xlabel={temps depuis le début de la séance (s)},ylabel={âge du cours affiché (s, échelle log)},
  grid=major,grid style={gray!20},tick label style={font=\small},label style={font=\small},table/col sep=comma,xmin=0,xmax=2450,
  ymin=0.01,ymax=5000,xtick={0,600,1200,1800,2400}]
\addplot[omThm,thick,mark=none] table[x=t,y=age_s]{figdata/platforms/23-market-data-distribution-and-entitlements/age_all.csv};
\addplot[omDef,thick,mark=none] table[x=t,y=age_s]{figdata/platforms/23-market-data-distribution-and-entitlements/age_conflate.csv};
\draw[gray,dashed] (axis cs:600,0.01) -- (axis cs:600,5000) node[pos=0.5,right,font=\scriptsize]{fin de séance};
\node[omThm,font=\scriptsize,anchor=north] at (axis cs:1500,250) {chaque message mis en file};
\node[omDef,font=\scriptsize,anchor=south west] at (axis cs:20,0.15) {en conflation par instrument};
\end{semilogyaxis}
\end{tikzpicture}

\omcaption{Âge des cours affichés par le tableau de bord de recherche, relevé toutes les 2\,000 livraisons (en file) ou toutes les 500 (en conflation). En file, l'âge
augmente tout au long de la séance, et l'arriéré met encore une demi-heure à se résorber~; en conflation, il reste voisin d'un dixième de seconde. Données~:
\texttt{fig\_entitle.py}.}\label{fig:pl:market-data-distribution-and-entitlements:age}
\end{omfigure}

\section{Licences, usages et unités de comptage}

Ce que la société paie n'est pas le flux, mais chacun de ses usages. Les plateformes distinguent l'usage en affichage, celui d'une personne qui regarde un écran, de
l'usage non-display par une machine, que le livre~1, chapitre~29, a introduit avec la redevance non-display~; elles les comptent chacun dans une unité déclarée.

\begin{definition}[Unité de comptage]\label{def:pl:market-data-distribution-and-entitlements:unit}
Une \emph{unité de comptage}\index{unité de comptage} est ce que la licence de la plateforme compte pour facturer ses données~: pour l'usage en affichage, typiquement chaque appareil ou
chaque identifiant utilisateur unique, de sorte qu'une personne qui dispose de deux identifiants compte deux fois~; pour l'usage non-display, typiquement chaque appareil (en général un
serveur) qui reçoit les données et en tire profit, ou chaque personne qui peut modifier l'application en temps réel.
\end{definition}

La définition est celle de la plateforme, et c'est là que les sociétés perdent de l'argent. La plateforme du chapitre --- Nasdaq, dont les documents figurent dans le registre --- compte un
abonné en affichage comme «~un appareil, un terminal informatique, un service automatisé, ou un couple identifiant utilisateur unique
et mot de passe~», et l'usage non-display par les serveurs qui reçoivent les données et en tirent profit, y compris ceux qui les calculent ou
en créent des \omterm{def:pl:market-data-distribution-and-entitlements:derived}{données dérivées} à partir d'elles. Les serveurs qui se contentent de transporter, de normaliser ou de stocker les données ne sont pas nécessairement assujettis à la redevance, mais la société doit pouvoir
montrer lesquels ils sont. Une société qui compte des personnes au lieu d'identifiants, ou des applications au lieu de serveurs, sous-compte par construction.

\section{Droits d'accès}

\begin{definition}[Droit d'accès, système de droits d'accès]\label{def:pl:market-data-distribution-and-entitlements:entitlement}
Un \emph{droit d'accès}\index{droit d'accès} (entitlement) est une autorisation accordée à un sujet --- un identifiant utilisateur ou une application --- de recevoir une source de
données pour un usage. Un \emph{système de droits d'accès}\index{système de droits d'accès} (entitlement system) détient les droits d'accès avec leurs dates d'octroi et de révocation,
répond à la question de savoir si un sujet a droit à un instant donné, et est vérifié par la plateforme de distribution avant toute livraison de données.
\end{definition}

La question est de savoir quand vérifier. La vérification à l'abonnement est peu coûteuse et suffit tant que rien ne change~; mais un trader qui quitte le desk, ou une
application dont la licence arrive à échéance, continue de recevoir jusqu'à son prochain abonnement. Dans la simulation du chapitre, le \omterm{def:pl:market-data-distribution-and-entitlements:entitlement}{droit d'accès} de l'utilisateur de l'écran est révoqué cinq
minutes après le début de la session. Si la vérification n'a lieu qu'à l'abonnement, l'écran reçoit encore 6\,005 mises à jour auxquelles il n'avait pas droit~; avec une vérification
à chaque livraison, il n'en reçoit aucune.

\omcode{code/firm/entitle/firm_entitle.py}{133}{147}{Livraison~: le droit d'accès est vérifié à nouveau à chaque livraison, de sorte qu'une révocation prenne
effet immédiatement, et l'âge de chaque message livré est enregistré.}\label{lst:pl:market-data-distribution-and-entitlements:serve}

Une plateforme de distribution réelle n'appelle pas le \omterm{def:pl:market-data-distribution-and-entitlements:entitlement}{système de droits d'accès} pour chacun des millions de messages~: elle met en cache la réponse par abonnement et invalide le
cache lorsqu'un \omterm{def:pl:market-data-distribution-and-entitlements:entitlement}{droit d'accès} change, ce qui produit le même effet pour un coût quasi nul. L'essentiel est que la révocation soit transmise
à la plateforme de distribution plutôt que découverte à la prochaine connexion.

\section{Rapports d'utilisation et audits des données de marché}

\begin{definition}[Rapport d'utilisation]\label{def:pl:market-data-distribution-and-entitlements:report}
Un \emph{rapport d'utilisation}\index{rapport d'utilisation} (usage report) est la déclaration périodique que la société adresse à une source de données pour indiquer combien d'\omterm{def:pl:market-data-distribution-and-entitlements:unit}{unités de comptage} ont utilisé ses données,
par produit et par usage, à partir de laquelle la source facture ses redevances.
\end{definition}

\begin{definition}[Audit des données de marché]\label{def:pl:market-data-distribution-and-entitlements:audit}
Un \emph{audit des données de marché}\index{audit des données de marché} (market-data audit) est l'examen, par une plateforme ou un fournisseur, de l'usage réel que fait une société de ses données sur une période écoulée,
rapproché de ce que la société a déclaré~; l'écart est refacturé, en général avec intérêts et, au-delà d'un seuil, les frais de l'audit.
\end{definition}

La société du chapitre a déclaré 40 abonnés en affichage chaque mois pendant trois ans, et aucun usage non-display. Son journal d'utilisation --- que la plateforme de distribution
enregistre à chaque abonnement --- raconte une tout autre histoire (la \cref{fig:pl:market-data-distribution-and-entitlements:usage}). Le nombre de traders est passé de 38
à 47, et six d'entre eux ont un second identifiant de connexion, si bien que les identifiants en affichage sont passés de 43 à 53~; les onze applications sont entrées
en service les unes après les autres et, à la fin, tournaient sur 22 serveurs.

\omcode{code/firm/entitle/firm_entitle.py}{168}{195}{Le rapport d'utilisation compte les unités de comptage distinctes par usage~; l'audit rapproche chaque mois
les licences détenues et valorise l'écart selon le barème des redevances.}\label{lst:pl:market-data-distribution-and-entitlements:audit}

\begin{omfigure}
\begin{tikzpicture}
\begin{axis}[width=11.5cm,height=5.6cm,xlabel={mois},ylabel={unités de comptage},
  grid=major,grid style={gray!20},tick label style={font=\small},label style={font=\small},table/col sep=comma,
  xmin=0,xmax=35,ymin=0,ymax=63,const plot]
\addplot[omDef,thick] table[x=month,y=display_used]{figdata/platforms/23-market-data-distribution-and-entitlements/usage.csv};
\addplot[omDef,thick,dashed] table[x=month,y=display_licensed]{figdata/platforms/23-market-data-distribution-and-entitlements/usage.csv};
\addplot[omThm,thick] table[x=month,y=nd_used]{figdata/platforms/23-market-data-distribution-and-entitlements/usage.csv};
\node[omDef,font=\scriptsize,anchor=south] at (axis cs:20,53) {identifiants en affichage utilisés};
\node[font=\scriptsize,anchor=north] at (axis cs:24,39) {abonnés en affichage déclarés~: 40};
\node[omThm,font=\scriptsize,anchor=south] at (axis cs:25,24) {serveurs non-display utilisés (déclarés~: 0)};
\end{axis}
\end{tikzpicture}

\omcaption{Trois ans du journal d'utilisation de la société modèle, comparés à ce qu'elle a déclaré. Les identifiants en affichage dépassent les 40 déclarés dès le premier mois
(43, à cause des seconds identifiants) et atteignent 53~; les serveurs non-display passent de 4 à 22 à mesure que les onze applications entrent en service. Données~:
\texttt{fig\_entitle.py}.}\label{fig:pl:market-data-distribution-and-entitlements:usage}
\end{omfigure}

Valorisé aux redevances 2026 de la plateforme pour son produit de profondeur --- 84~\$ par mois et par abonné professionnel en affichage, et pour l'usage non-display 412~\$ par
mois par abonné jusqu'à 39, puis un forfait de 16\,490~\$ par société de 40 à 99 (la \cref{fig:pl:market-data-distribution-and-entitlements:fees}) ---
l'écart est de 1\,900~\$ le premier mois et de 10\,156~\$ le dernier, soit 258\,596~\$ sur les trois ans~: 22\,932~\$ pour l'affichage et 235\,664~\$ pour
non-display. Les unités non déclarées atteignent 87.5~\% de celles qui ont été déclarées, très au-delà des 10~\% à partir desquels le contrat d'abonnement de la plateforme ajoute les
frais de l'audit~; les intérêts s'ajoutent, à un taux que le modèle ne valorise pas. L'écart en affichage est la moindre part. La vraie erreur de la société
a été de croire que des données qu'elle avait payées pour les recevoir pouvaient être utilisées par n'importe quel programme.

\begin{omfigure}
\begin{tikzpicture}
\begin{axis}[width=11.5cm,height=5.4cm,xlabel={abonnés non-display (serveurs)},ylabel={redevance mensuelle (milliers de \$)},
  grid=major,grid style={gray!20},tick label style={font=\small},label style={font=\small},table/col sep=comma,
  xmin=0,xmax=120,ymin=0,ymax=38,const plot]
\addplot[omThm,thick] table[x=devices,y expr=\thisrow{fee}/1000]{figdata/platforms/23-market-data-distribution-and-entitlements/fees.csv};
\draw[omDef,dashed] (axis cs:22,0) -- (axis cs:22,9.064);
\draw[omDef,-{Stealth[length=1.5mm]}] (axis cs:55,6) node[right,font=\scriptsize]{la société modèle~: 22 serveurs, 9\,064~\$} -- (axis cs:23,8.8);
\node[font=\scriptsize,anchor=south east] at (axis cs:37,17.5) {39~: 16\,068~\$};
\node[font=\scriptsize,anchor=south] at (axis cs:70,17.5) {40 à 99~: 16\,490~\$};
\node[font=\scriptsize,anchor=north east] at (axis cs:98,32) {100 à 249~: 32\,990~\$};
\end{axis}
\end{tikzpicture}

\omcaption{La redevance mensuelle non-display 2026 de la plateforme pour son produit de profondeur, en fonction du nombre d'abonnés assujettis à la redevance~: 412~\$ chacun jusqu'à 39,
puis forfaitaire par société. Entre 40 et 99 serveurs, la redevance ne change pas, et c'est pourquoi le compte et le palier, et non le prix, déterminent la facture.
Données~: \texttt{fig\_entitle.py}.}\label{fig:pl:market-data-distribution-and-entitlements:fees}
\end{omfigure}

Le remède est aussi technique que contractuel. Le journal d'utilisation de la plateforme de distribution est la preuve qu'exige un audit, et il doit enregistrer l'\omterm{def:pl:market-data-distribution-and-entitlements:unit}{unité de
comptage} qu'utilise la plateforme --- des identifiants et non des personnes~; des serveurs et non des applications --- pendant toute la période de rétrospection. La plateforme du chapitre limite
sa responsabilité pour une erreur de bonne foi aux trois années précédant la découverte de la sous-déclaration, si bien que trois ans de journaux constituent le minimum. Un
\omterm{def:pl:post-trade-systems:recon}{rapprochement} mensuel du journal avec les licences, mené par la société avant qu'un auditeur ne le fasse, transforme un rappel de facturation en achat de licence.

\section{Données dérivées}

\begin{definition}[Données dérivées]\label{def:pl:market-data-distribution-and-entitlements:derived}
Les \emph{données dérivées}\index{données dérivées} (derived data) sont des informations calculées en tout ou partie à partir des données de la plateforme, et qui ne permettent
ni de les reconstituer ni de s'y substituer raisonnablement~: un indice, un signal, la fair value d'un modèle. Ce qui relève ou non des données dérivées est fixé par
la politique de chaque plateforme.
\end{definition}

Une cotation conflatée n'est pas une donnée dérivée~: c'est la donnée de la plateforme, livrée moins souvent, et elle requiert la même licence. Un prix milieu publié chaque
seconde peut être un substitut de la cotation~; une estimation de volatilité ou la fair value d'un panier ne l'est généralement pas. La distinction compte dans deux
sens~: les \omterm{def:pl:market-data-distribution-and-entitlements:derived}{données dérivées} peuvent être redistribuées selon des conditions que les données brutes n'autorisent pas, et les serveurs qui les créent sont, pour la
plateforme du chapitre, parmi les \omterm{def:pl:market-data-distribution-and-entitlements:unit}{unités de comptage} non-display. Le modèle de \omterm{def:pl:market-data-distribution-and-entitlements:entitlement}{droits d'accès} doit donc porter l'usage --- affichage, non-display,
dérivé --- et non pas seulement la source, et la \omterm{def:pl:a-tick-store-on-open-formats:store}{base de ticks} du chapitre~\ref{ch:pl:a-tick-store-on-open-formats} doit enregistrer sous quelle licence
chaque jeu de données a été capturé, puisque la licence suit les données jusque dans la recherche.

\begin{dated}{2026-09}{Politiques de données et redevances d'une plateforme}\label{dat:pl:market-data-distribution-and-entitlements:nasdaq}
La grille tarifaire des actions américaines de Nasdaq fixe, à compter du 1\up{er} janvier 2026, les redevances mensuelles de Nasdaq TotalView à 84.00~\$ par abonné professionnel (80.50~\$ en 2025) et à 15.00~\$ par abonné non professionnel~;
pour l'usage non-display de ses données de profondeur avec accès direct, 412.00~\$ par abonné de 1 à 39,
16\,490~\$ par société de 40 à 99, 32\,990~\$ de 100 à 249 et 75\,000~\$ à partir de 250~; et des redevances de distributeur de 1\,690~\$ par société pour la distribution
interne et 3\,340~\$ pour l'accès direct. Ses US Equities and Options Data Policies (version 2.6) définissent l'abonné, l'usage non-display et les
\omterm{def:pl:market-data-distribution-and-entitlements:derived}{données dérivées} tels que ce chapitre les définit, et précisent qu'un usage non déclaré d'un système ne peut pas être compensé sur la période d'audit.
Son contrat d'abonné limite la responsabilité pour une erreur de bonne foi à trois ans de redevances impayées avec intérêts, et ajoute les frais de l'audit lorsque la sous-déclaration
atteint 10~\% des unités déclarées.
\end{dated}

\section{Tutoriel~: distribuer, octroyer des droits, rapporter, auditer}\label{tut:pl:market-data-distribution-and-entitlements:tut}

\begin{tutorial}
\textbf{Objectif.} Distribuer un flux aux abonnés de vitesses différentes, faire respecter les \omterm{def:pl:market-data-distribution-and-entitlements:entitlement}{droits d'accès} et réconcilier trois ans d'usage avec les
licences détenues. \textbf{État final~:} le \cref{tab:pl:market-data-distribution-and-entitlements:policies} et le montant de l'audit
258\,596~\$.
\begin{tutsteps}
\item \textbf{Flux}~: \texttt{pl\_entitle.updates()}, dix minutes sur vingt instruments.
\item \textbf{Distribution}~: \texttt{distribute(ups, \textquotedbl all\textquotedbl)} et
\texttt{distribute(ups, \textquotedbl conflate\textquotedbl)}~; comparer les files d'attente et les délais.
\item \textbf{\omterm{def:pl:market-data-distribution-and-entitlements:entitlement}{Droits d'accès}}~: relancer avec \texttt{check\_on\_delivery=False} et compter les livraisons à l'écran après la révocation.
\item \textbf{Usage}~: \texttt{usage\_log()} et \texttt{firm\_entitle.usage\_report} pour un mois quelconque.
\item \textbf{Audit}~: \texttt{audit()}, les mois, les insuffisances de déclaration et le montant dû.
\end{tutsteps}
\textbf{Pour aller plus loin.} Déplacer le tableau de bord de recherche sur les serveurs d'un niveau de distribution et voir ce que devient le décompte non-display~;
chiffrer l'audit avec les redevances de chaque année au lieu de celles de 2026.
\end{tutorial}

\section{Construction~: droits d'accès}\label{bld:pl:market-data-distribution-and-entitlements:entitle}

\begin{build}
\bfield{Rôle} Livrer les données de marché à l'intérieur de la société à chaque abonné sous la forme qu'il peut utiliser, uniquement à ceux qui y ont droit, avec la preuve de
chaque usage comptabilisé comme la plateforme le compte.
\bfield{Interface} \texttt{Entitlement}, \texttt{EntitlementSystem} (\texttt{grant}, \texttt{revoke}, \texttt{check}), \texttt{Bus}
(\texttt{subscribe}, \texttt{run}), \texttt{UsageLog}, \texttt{usage\_report}, \texttt{audit}.
\bfield{Règles} Une file par abonné~; conflation pour les abonnés qui veulent l'état~; \omterm{def:pl:market-data-distribution-and-entitlements:entitlement}{droit d'accès} vérifié à l'abonnement et appliqué à la
livraison~; chaque usage consigné dans l'\omterm{def:pl:market-data-distribution-and-entitlements:unit}{unité de comptage} de la plateforme~; journaux conservés au moins aussi longtemps que la période de rétrospection~; une réconciliation mensuelle avant tout
audit.
\bfield{Tests d'acceptation} \texttt{code/firm/entitle/tests/}~: vérifications par sujet, source, usage et horodatage~; livraison en file et conflatée~; un abonné très
sollicité, ordonnancé correctement~; révocation vérifiée à la livraison plutôt qu'à l'abonnement~; décomptes d'usage par identifiant et par appareil, et
l'insuffisance et le montant de l'audit.
\bfield{Pour aller plus loin} Redevances par paliers avec licences d'entreprise~; un registre des \omterm{def:pl:market-data-distribution-and-entitlements:derived}{données dérivées}~; un \omterm{def:pl:market-data-distribution-and-entitlements:report}{rapport d'utilisation} au format de fichier d'une plateforme.
\end{build}

\begin{omsources}
\item Nasdaq, \emph{U.S. Equities Price List}, redevances en vigueur en 2025, 2026 et 2027.
\item Nasdaq, \emph{US Equities and Options Data Policies}, version 2.6.
\item Nasdaq, \emph{Global Subscriber Agreement}, sections 5(g) et 5(h).
\item Le livre~1, chapitres~9 et~29 (flux directs, redevances non-display)~; le livre~13, chapitre~12 (consommateurs lents et conflation).
\end{omsources}

\section{Exercices}

\begin{exercise}[$\star$]\label{exo:pl:market-data-distribution-and-entitlements:1}
Expliquer pourquoi la plateforme garde une file par abonné plutôt qu'une file par sujet.
\end{exercise}

\begin{exercise}[$\star$]\label{exo:pl:market-data-distribution-and-entitlements:2}
Un trader utilise les mêmes données de profondeur sur un poste de travail et, avec un second identifiant de connexion, sur un ordinateur portable. Calculer le nombre d'abonnés à l'affichage que cela représente pour la
plateforme du chapitre.
\end{exercise}

\begin{exercise}[$\star$]\label{exo:pl:market-data-distribution-and-entitlements:3}
Calculer la redevance non-display de la société modèle pour le dernier mois, et ce qu'elle deviendrait avec 45 serveurs.
\end{exercise}

\begin{exercise}[$\star\star$]\label{exo:pl:market-data-distribution-and-entitlements:4}
Déterminer lequel des quatre abonnés doit être conflaté, et expliquer pourquoi ce n'est pas le cas de la stratégie.
\end{exercise}

\begin{exercise}[$\star\star$]\label{exo:pl:market-data-distribution-and-entitlements:5}
Expliquer pourquoi le tableau de bord de recherche, alimenté par une file d'attente, affiche en moyenne des cours vieux de quinze minutes alors qu'il met 10~millisecondes par message.
\end{exercise}

\begin{exercise}[$\star\star$]\label{exo:pl:market-data-distribution-and-entitlements:6}
La société veut vérifier les \omterm{def:pl:market-data-distribution-and-entitlements:entitlement}{droits d'accès} à chaque livraison sans appeler le \omterm{def:pl:market-data-distribution-and-entitlements:entitlement}{système de droits d'accès} des millions de fois par minute. Expliquer comment s'y prendre.
\end{exercise}

\begin{exercise}[$\star\star\star$]\label{exo:pl:market-data-distribution-and-entitlements:7}
\emph{Programmation.} Relancer \texttt{audit} avec la société qui déclare correctement ses identifiants d'affichage chaque mois mais sans aucun usage non-display. Calculer le
montant dû.
\end{exercise}

\begin{exercise}[$\star\star\star$]\label{exo:pl:market-data-distribution-and-entitlements:8}
\emph{Trouver la faille.} «~Nous payons la redevance de distributeur pour la distribution interne, donc toute application interne à la société peut utiliser les
données.~»
\end{exercise}

\section{Problème~: onze applications}

\begin{problem}[{Problème du week-end --- onze applications}]\label{pb:pl:market-data-distribution-and-entitlements:1}
La plateforme du chapitre, ses quatre abonnés et les trois années d'usage de la société modèle.

\medskip\noindent\textbf{Partie I --- Distribution.}
\begin{enumerate}
\item Que fait la plateforme de distribution que ne fait pas le raccordement de chaque consommateur directement au flux~?
\item À quelle vitesse chaque abonné traite-t-il les messages, et à quelle vitesse le flux arrive-t-il~?
\item Que devient le tableau de bord de recherche lorsque chaque message est mis en file d'attente~?
\item Que change la conflation pour le service de risque et pour le tableau de bord~?
\item Quelle est la plus grande file de la stratégie, et pourquoi~?
\end{enumerate}

\medskip\noindent\textbf{Partie II --- \omterm{def:pl:market-data-distribution-and-entitlements:entitlement}{Droits d'accès}.}
\begin{enumerate}[resume]
\item Que désigne un \omterm{def:pl:market-data-distribution-and-entitlements:entitlement}{droit d'accès}~?
\item Que se passe-t-il après une révocation lorsque les \omterm{def:pl:market-data-distribution-and-entitlements:entitlement}{droits d'accès} ne sont vérifiés qu'à l'abonnement~?
\item Comment une plateforme applique-t-elle à moindre coût les \omterm{def:pl:market-data-distribution-and-entitlements:entitlement}{droits d'accès} à chaque livraison~?
\item Quelle est l'\omterm{def:pl:market-data-distribution-and-entitlements:unit}{unité de comptage} de l'usage à l'affichage et de l'usage non-display sur la plateforme du chapitre~?
\item Quels serveurs ne sont pas nécessairement assujettis à redevance~?
\end{enumerate}

\medskip\noindent\textbf{Partie III --- L'audit.}
\begin{enumerate}[resume]
\item Qu'a déclaré la société, et que montrent son journal d'usage pour le premier et le dernier mois~?
\item Quels sont les frais mensuels utilisés pour chiffrer l'écart~?
\item Que doit-on pour le premier et le dernier mois~?
\item Que doit-on sur les trois années, pour l'affichage et pour le non-display~?
\item Pourquoi l'audit coûte-t-il plus cher à la société que les frais rétroactifs~?
\end{enumerate}

\medskip\noindent\textbf{Partie IV --- Le verdict.}
\begin{enumerate}[resume]
\item Énoncer le \emph{résultat nommé}~: l'écart de licence pour les utilisateurs en affichage et les applications en non-display, mis en évidence par le \omterm{def:pl:post-trade-systems:recon}{rapprochement} de trois années de
journaux d'usage, et son coût rétroactif aux tarifs 2026 de la plateforme.
\item Pendant combien de temps les journaux d'usage doivent-ils être conservés, et dans quelle unité~?
\item Une cotation conflatée est-elle une donnée dérivée~?
\item Quelle modification faudrait-il apporter en premier à la plateforme de la société~?
\item En une phrase~: à quoi sert une licence de données de marché~?
\end{enumerate}
\end{problem}

\section{Questions d'entretien}

\begin{interviewq}[$\star$ \iqroles{développeur}]\label{iq:pl:market-data-distribution-and-entitlements:1}
Qu'est-ce que la conflation, et quand l'utiliseriez-vous~?
\end{interviewq}

\begin{interviewq}[$\star\star$ \iqroles{développeur}]\label{iq:pl:market-data-distribution-and-entitlements:2}
Un consommateur lent bloque votre bus de données de marché. Que faites-vous~?
\end{interviewq}

\begin{interviewq}[$\star\star$ \iqroles{développeur}]\label{iq:pl:market-data-distribution-and-entitlements:3}
Quelle est la différence entre l'usage en affichage et l'usage non-display des données de marché, et pourquoi cela importe-t-il à un développeur~?
\end{interviewq}

\begin{interviewq}[$\star\star$ \iqroles{développeur}]\label{iq:pl:market-data-distribution-and-entitlements:4}
Une bourse annonce un audit de votre usage des données sur les trois dernières années. De quoi devez-vous disposer~?
\end{interviewq}

\begin{interviewq}[$\star\star\star$ \iqroles{développeur}]\label{iq:pl:market-data-distribution-and-entitlements:5}
Concevez la \omterm{def:pl:market-data-distribution-and-entitlements:platform}{plateforme de distribution des données de marché} et des \omterm{def:pl:market-data-distribution-and-entitlements:entitlement}{droits d'accès} d'une société comptant mille utilisateurs et cent applications.
\end{interviewq}
